% =====================================================================
%  Section 3 -- Estimands (owner: THEORY). Proofs in appendix/A-proofs-estimands.tex.
% =====================================================================
\providecommand{\osc}{\operatorname{osc}}
\providecommand{\dTV}{d_{\mathrm{TV}}}

\section{Estimands: what does it mean to amplify?}\label{sec:estimands}

To say that a recommender \emph{amplifies} antisemitic content is to make a comparison. Exposure under the algorithm is said to be larger than exposure under something else, for someone. The claim has no truth value until both the ``something else'' and the ``someone'' are named, and different choices answer different questions. This section names them. \Cref{sec:est-setup} defines exposure as a potential outcome indexed by a policy and a demand profile. \Cref{def:amplification} defines amplification at fixed demand, relative to a named reference policy, and \cref{sec:est-reference} argues that the reference is part of the estimand rather than a nuisance. \Cref{sec:est-decomp} turns to real users, whose demand responds to what they are shown. Their observed exposure splits into a demand-driven baseline, a supply effect, a demand-shaping effect and an interaction between the last two (\cref{thm:decomp}). \Cref{sec:est-audit} states what an audit with sock puppets of fixed demand can identify (\cref{prop:audit-scope}): the supply effect at the audited personas, and nothing about how the algorithm shapes demand.

The framework is the potential-outcomes framework of \citet{rubin1974} and \citet{imbensrubin2015}. The amplification ratio has a precedent in \citet{huszar2022}, who compared users randomly assigned to an algorithmically ranked home timeline with a control group that kept a reverse-chronological one. Feedback between recommendations and preferences has been studied by simulation \citep{chaney2018}, as a dynamical system \citep{jiang2019}, and by estimating the preference shifts that a recommender induces \citep{carroll2022}. The two-factor decomposition is a mediation analysis in which demand is the mediator \citep{robins1992,pearl2001,vanderweele2014,vanderweele2014b}. Its symmetric version is the decomposition of a difference between two rates into a rate and a composition component \citep{kitagawa1955}, which is also the Shapley value \citep{shapley1953} of a two-player game. What is new is the application to amplification audits: the bilinear representation \eqref{eq:est-bilinear}, the characterisation and bounds of the interaction in \cref{thm:decomp}(d), and the precise statement in \cref{prop:audit-scope} of what fixed-demand audits identify. \Cref{app:proofs-estimands} gives the proofs.

\subsection{Policies, demand profiles and exposure}\label{sec:est-setup}

An \emph{account} $u$ is shown content in \emph{sessions}. In a session it is exposed to a sequence of items $c_1,\dots,c_T$. Whether ``exposed'' means an impression or a consumption is fixed by the protocol (\cref{sec:protocol}); the theory is the same. Each item $c$ has a true label $L_c\in\{0,1\}$, with $L_c=1$ for antisemitic content. Classifier labels and their errors enter in \cref{sec:measurement}. The \emph{exposure} of $u$ in a session is the share of antisemitic items,
\begin{equation}\label{eq:est-exposure}
  E_u=\frac1T\sum_{t=1}^{T}\ind{L_{c_t}=1}.
\end{equation}

Two rules determine which items an account sees. The \emph{policy} $\pi$ is the platform's rule for selecting and ordering a slate of items, given the account's history. The \emph{demand profile} $d\in\mathcal D$ is the account's rule for acting on what it is shown: which accounts it follows, what it searches for, which items it selects from a slate and how long it dwells on them. For a sock puppet, $d$ is a script fixed by the researcher, and we call such a fixed demand profile a \emph{persona}. For a real user, $d$ summarises preferences and habits at a given time, and it can change. Items are drawn from an \emph{eligible pool} $\mathcal C_t$: the items that exist at time $t$, have survived moderation and may be shown to the account. The \emph{algorithmic policy} $\mathrm{alg}$ is the platform's production recommender. A \emph{reference policy} $\rho$ is a policy with which it is compared. Three references recur in the audit literature and in this paper:
\begin{itemize}
  \item \emph{reverse-chronological} ($\mathrm{chron}$): items posted by the accounts that $d$ follows, newest first;
  \item \emph{random} ($\mathrm{rand}$): items drawn uniformly at random from $\mathcal C_t$, or from a declared sub-pool such as the items in the account's language;
  \item \emph{popularity-only} ($\mathrm{pop}$): items from $\mathcal C_t$ ranked by an engagement score that is not personalised.
\end{itemize}
Let $\Pi=\{\mathrm{alg}\}\cup\Pi_{\mathrm{ref}}$, where $\Pi_{\mathrm{ref}}$ is the set of references in use. Let $E_u(\pi;d)$ be the exposure of account $u$ in a session served by policy $\pi$ while the account acts according to $d$. For a puppet with persona $d_u$, the exposure $E_u(\pi)$ of the notation table is $E_u(\pi;d_u)$. The \emph{mean exposure}
\begin{equation}\label{eq:est-mu}
  \mu_\pi(d)=\E\{E_u(\pi;d)\},\qquad \pi\in\Pi,\ d\in\mathcal D,
\end{equation}
averages over everything that is random given $(\pi,d)$: the policy's own randomisation, the contents of the pool during the audit window, the account's stochastic choices and the state in which a session starts. Writing $\mu_\pi$ as a function of $d$ alone presumes that accounts do not interfere with one another and that the policy is stable over the audit window. \Cref{sec:identification} states these conditions and relaxes them, and \cref{sec:markov} represents $\mu_\pi(d)$ through a Markov chain.

\subsection{Amplification at fixed demand}\label{sec:est-amp}

\begin{definition}[Amplification]\label{def:amplification}
Fix a reference policy $\rho\in\Pi_{\mathrm{ref}}$. For a persona $d$ with $\mu_\rho(d)>0$, the \emph{amplification ratio} and the \emph{additive amplification} of antisemitic content by the algorithmic policy, relative to $\rho$, are
\[
  \mathrm{AR}_\rho(d)=\frac{\mu_{\mathrm{alg}}(d)}{\mu_\rho(d)},\qquad
  \mathrm{AA}_\rho(d)=\mu_{\mathrm{alg}}(d)-\mu_\rho(d).
\]
For a probability distribution $Q$ on $\mathcal D$, write $\mu_\pi(Q)=\int\mu_\pi\,dQ$. The \emph{population-weighted} amplification ratio and additive amplification are
\[
  \mathrm{AR}_\rho(Q)=\frac{\mu_{\mathrm{alg}}(Q)}{\mu_\rho(Q)},\qquad
  \mathrm{AA}_\rho(Q)=\mu_{\mathrm{alg}}(Q)-\mu_\rho(Q)=\int\mathrm{AA}_\rho\,dQ .
\]
The algorithm \emph{amplifies} relative to $\rho$ (for $d$, or for $Q$) if $\mathrm{AR}_\rho>1$, equivalently $\mathrm{AA}_\rho>0$, and \emph{de-amplifies} if $\mathrm{AR}_\rho<1$.
\end{definition}

The distribution $Q$ is either the persona mix of an audit, given by its design weights, or the distribution of a target population over $\mathcal D$. When a single reference is in play we write $\mathrm{ref}$ for it and drop the subscript. $\mathrm{AA}$ is always defined, and it is the scale on which exposure burdens add up across users. $\mathrm{AR}$ is the natural scale when exposure is rare, and it is undefined when the reference delivers no antisemitic content at all. The population ratio $\mathrm{AR}_\rho(Q)$ is a ratio of means, not a mean of ratios. The difference matters, because personas differ in baseline exposure.

\begin{lemma}[Aggregation and change of reference]\label{lem:est-agg}
Let $\mu_\rho>0$ on a set of $Q$-probability one.
\begin{enumerate}[(a)]
  \item $\mathrm{AR}_\rho(Q)=\int\mathrm{AR}_\rho\,w_Q\,dQ$, with weights $w_Q=\mu_\rho/\mu_\rho(Q)$ that integrate to one.
  \item If $\int\mathrm{AR}_\rho\,dQ<\infty$, then $\int\mathrm{AR}_\rho\,dQ-\mathrm{AR}_\rho(Q)=-\Cov_Q(\mathrm{AR}_\rho,\mu_\rho)/\mu_\rho(Q)$.
  \item For two references, $\mathrm{AR}_\rho(d)=\mathrm{AR}_{\rho'}(d)\,\mu_{\rho'}(d)/\mu_\rho(d)$. For a ladder of policies $\rho_0,\rho_1,\dots,\rho_k=\mathrm{alg}$ with positive means, $\mathrm{AR}_{\rho_0}(d)=\prod_{j=1}^{k}\mu_{\rho_j}(d)/\mu_{\rho_{j-1}}(d)$.
\end{enumerate}
\end{lemma}

\begin{proof}
(a) Integrate $\mu_{\mathrm{alg}}=\mathrm{AR}_\rho\,\mu_\rho$ against $Q$ and divide by $\mu_\rho(Q)$. (b) By (a), $\mathrm{AR}_\rho(Q)=\E_Q[\mathrm{AR}_\rho\,\mu_\rho]/\E_Q[\mu_\rho]=\E_Q[\mathrm{AR}_\rho]+\Cov_Q(\mathrm{AR}_\rho,\mu_\rho)/\E_Q[\mu_\rho]$. (c) The products telescope.
\end{proof}

Part (a) says that the population ratio weights each persona's ratio by its exposure under the reference. A persona's contribution to $\mathrm{AR}_\rho(Q)$ is $\mathrm{AR}_\rho(d)\,w_Q(d)\,Q(d)=Q(d)\,\mu_{\mathrm{alg}}(d)/\mu_\rho(Q)$, its share of exposure under the algorithm, which does not depend on $\mu_\rho(d)$. An enormous persona ratio produced by a tiny denominator therefore cannot make $\mathrm{AR}_\rho(Q)$ enormous by itself. But personas with almost no reference exposure can still dominate $\mathrm{AR}_\rho(Q)$ if they account for most of the exposure under the algorithm. Part (b) says that the unweighted average of persona ratios overstates the population ratio when the personas with the lowest reference exposure have the largest ratios. That is the pattern to expect if an algorithm leads users to content that they would not have found by themselves. We therefore report $\mathrm{AR}_\rho(Q)$ and $\mathrm{AA}_\rho(Q)$ for a declared $Q$, report persona-level values separately, and never present an average of persona ratios as a population quantity.

\subsection{The reference policy is part of the estimand}\label{sec:est-reference}

By \cref{lem:est-agg}(c), amplification relative to $\rho$ and amplification relative to $\rho'$ differ by the factor $\mu_{\rho'}(d)/\mu_\rho(d)$, which does not involve the algorithm. The same algorithm can therefore amplify relative to one reference and de-amplify relative to another. There is no reference-free notion of amplification, and the choice of reference is a choice of question.
\begin{itemize}
  \item $\mathrm{AR}_{\mathrm{chron}}$ asks how ranking, including the recommendation of items from accounts that the persona does not follow, changes exposure relative to the persona's own subscriptions. The follow graph is part of $d$ and is held fixed. If $d$ follows no source of antisemitic content, then $\mu_{\mathrm{chron}}(d)=0$, the ratio is undefined, and $\mathrm{AA}_{\mathrm{chron}}(d)=\mu_{\mathrm{alg}}(d)$: all of the persona's exposure is attributed to the algorithm.
  \item $\mathrm{AR}_{\mathrm{rand}}$ asks how selective the algorithm is relative to the base rate in the eligible pool. If exposure counts impressions, $\mu_{\mathrm{rand}}(d)$ is the time-averaged prevalence of antisemitic content in the pool, which is the same for all personas that share a pool, and $\mathrm{AR}_{\mathrm{rand}}$ is a selection ratio. If exposure counts consumption, $\mu_{\mathrm{rand}}(d)$ also reflects the persona's choices from random slates.
  \item $\mathrm{AR}_{\mathrm{pop}}$ asks what personalisation adds to a ranking by global engagement.
\end{itemize}
Applied to the ladder $\mathrm{rand}\to\mathrm{pop}\to\mathrm{alg}$, \cref{lem:est-agg}(c) gives $\mathrm{AR}_{\mathrm{rand}}(d)=\{\mu_{\mathrm{pop}}(d)/\mu_{\mathrm{rand}}(d)\}\,\mathrm{AR}_{\mathrm{pop}}(d)$: selection by popularity times personalisation. A ladder of references thus decomposes amplification by mechanism, exactly and multiplicatively. The decomposition depends on which rungs are chosen, but for a given ladder there is no ordering ambiguity of the kind discussed in \cref{sec:est-decomp}.

Three further choices belong to the estimand.
\begin{enumerate}[(i)]
  \item \emph{The pool.} Every policy in \cref{def:amplification} draws from the same pool $\mathcal C_t$. Moderation that removes items before ranking changes $\mu_\pi$ for every $\pi$, and it is not amplification. ``Supply'' therefore has two layers: the \emph{content supply}, meaning what producers create and moderation leaves eligible, and the \emph{algorithmic supply}, meaning what the policy selects from it. Amplification concerns the second. Producers who respond to what the algorithm rewards change the first. That slower feedback, central to the supply-and-demand account of \citet{munger2022}, lies outside every estimand in this paper.
  \item \emph{The normalisation.} $E_u$ is a share of a fixed number $T$ of items. If the algorithm also lengthens sessions, total exposure changes by the factor $T_{\mathrm{alg}}/T_{\mathrm{ref}}$ in addition to $\mathrm{AR}$. We report shares, and session lengths separately.
  \item \emph{Implementability.} A reference should be a policy that can be run. Reverse-chronological and popularity-only feeds can sometimes be emulated from public data, whereas a random draw from the eligible pool requires access to the pool, which in practice means platform-side data access (\cref{sec:protocol}). A reference that cannot be implemented can still be defined, but amplification relative to it is then identified only under a model.
\end{enumerate}

\subsection{Dynamic demand: supply, demand shaping and feedback}\label{sec:est-decomp}

A puppet's demand is fixed by design. A real user's is not. What a user is shown today can change what the user seeks tomorrow \citep{jiang2019,carroll2022}, and a recommender trained on behaviour that it has itself shaped feeds the change back into its recommendations \citep{chaney2018}. The observed exposure of real users then reflects three things: what they would have sought anyway, what the algorithm puts in front of them, and how the algorithm has changed what they seek.

We formalise this with a \emph{policy schedule}. A formation period $W_0$ precedes an evaluation window $W$. Policy $b$ serves the user during $W_0$, and policy $a$ during $W$. Let $E_u(a;b)$ be the user's exposure in $W$, with $E_u(a)=E_u(a;a)$. For a population $\mathcal U$ of real users and a user $U$ drawn from it, with population weights if needed, write
\[
  \bar\mu(a;b)=\E\{E_U(a;b)\},\qquad a,b\in\{\mathrm{alg},\mathrm{ref}\}.
\]
Each of the four cells is the outcome of a feasible intervention on the platform's policy schedule. None is a cross-world counterfactual. The diagonal cells are the mean exposures in a world run by the algorithm throughout, which is what an observer of the current platform sees, and in a world run by the reference throughout. The off-diagonal cells switch the policy at the start of $W$. Thus $\bar\mu(\mathrm{ref};\mathrm{alg})$ is what the reference would deliver to users whose demand the algorithm has formed, and $\bar\mu(\mathrm{alg};\mathrm{ref})$ is what the algorithm would deliver to users whose demand formed under the reference. Let $D_U(b)\in\mathcal D$ be the user's demand profile at the start of $W$ after formation under $b$, and let $F_b$ be its distribution over the population.

\begin{assumption}[Demand sufficiency]\label{asm:demand}
For $a,b\in\{\mathrm{alg},\mathrm{ref}\}$, $\E\{E_U(a;b)\mid D_U(b)=d\}=\mu_a(d)$ for $F_b$-almost every $d$, where $\mu_a$ is the persona-level mean exposure \eqref{eq:est-mu}.
\end{assumption}

\Cref{asm:demand} has two parts, and it presumes that demand does not change materially during $W$, which is short relative to $W_0$: $\mu_a(d)$ is the exposure of a persona whose demand stays at $d$. First, the formation regime affects exposure in $W$ only through the demand profile. This fails if the policy carries per-user state, such as a learned embedding, that depends on $b$ beyond what $d$ records. That state is then a third factor, and the two-player game below becomes a three-player game with the same Shapley logic. Second, a user with demand profile $d$ is, as far as the policy is concerned, exchangeable with a puppet of persona $d$. This is the transport assumption behind every sock-puppet audit, and \cref{sec:identification} returns to it. Under \cref{asm:demand}, iterated expectations give
\begin{equation}\label{eq:est-bilinear}
  \bar\mu(a;b)=\int\mu_a\,dF_b=\mu_a(F_b),
\end{equation}
which is bilinear in the \emph{supply response} $\mu_a$ and the \emph{demand distribution} $F_b$.

Define the total algorithmic effect and its components:
\begin{equation}\label{eq:est-components}
\begin{aligned}
  \tau&=\bar\mu(\mathrm{alg};\mathrm{alg})-\bar\mu(\mathrm{ref};\mathrm{ref}), &
  \tau_{\mathrm{sup}}&=\bar\mu(\mathrm{alg};\mathrm{ref})-\bar\mu(\mathrm{ref};\mathrm{ref}),\\
  \iota&=\bar\mu(\mathrm{alg};\mathrm{alg})-\bar\mu(\mathrm{alg};\mathrm{ref})-\bar\mu(\mathrm{ref};\mathrm{alg})+\bar\mu(\mathrm{ref};\mathrm{ref}), &
  \tau_{\mathrm{dem}}&=\bar\mu(\mathrm{ref};\mathrm{alg})-\bar\mu(\mathrm{ref};\mathrm{ref}).
\end{aligned}
\end{equation}
The \emph{supply effect} $\tau_{\mathrm{sup}}$ is the change in exposure from switching the policy while demand stays as the reference formed it. The \emph{demand-shaping effect} $\tau_{\mathrm{dem}}$ is the change caused by the algorithm's effect on demand, seen through the reference. The \emph{feedback interaction} $\iota$ is the rest. On the ratio scale, write $\mathrm{TR}=\bar\mu(\mathrm{alg};\mathrm{alg})/\bar\mu(\mathrm{ref};\mathrm{ref})$ for the total ratio, $\mathrm{AR}(F_b)=\bar\mu(\mathrm{alg};b)/\bar\mu(\mathrm{ref};b)$ for the amplification ratio at the demand formed under $b$, and $\mathrm{DR}(a)=\bar\mu(a;\mathrm{alg})/\bar\mu(a;\mathrm{ref})$ for the demand ratio under policy $a$. Under \cref{asm:demand}, $\mathrm{AR}(F_b)=\mathrm{AR}_{\mathrm{ref}}(F_b)$ in the sense of \cref{def:amplification}.

The attribution problem is a two-player game. The players are \emph{supply} and \emph{demand}. For a coalition $J\subseteq\{\mathrm{sup},\mathrm{dem}\}$ let $a_J=\mathrm{alg}$ if $\mathrm{sup}\in J$ and $a_J=\mathrm{ref}$ otherwise, define $b_J$ in the same way from $\mathrm{dem}$, and let $v(J)=\bar\mu(a_J;b_J)-\bar\mu(\mathrm{ref};\mathrm{ref})$. Let $\osc_S f=\sup_Sf-\inf_Sf$, let $\dTV$ be the total-variation distance, and let $W_1$ be the Wasserstein-1 distance.

\begin{theorem}[Supply/demand decomposition]\label{thm:decomp}
\leavevmode
\begin{enumerate}[(a)]
  \item \emph{Two paths.} Without further assumptions, $\bar\mu(\mathrm{alg};\mathrm{alg})=\bar\mu(\mathrm{ref};\mathrm{ref})+\tau_{\mathrm{sup}}+\tau_{\mathrm{dem}}+\iota$, and
  \[
    \tau=\tau_{\mathrm{sup}}+(\tau_{\mathrm{dem}}+\iota)=(\tau_{\mathrm{sup}}+\iota)+\tau_{\mathrm{dem}}.
  \]
  The supply contribution is $\tau_{\mathrm{sup}}$ if the policy is switched first, at demand formed by the reference, and $\tau_{\mathrm{sup}}+\iota$ if demand is switched first, so that the policy acts on demand formed by the algorithm. The two path attributions differ by $\iota$, and so do the two demand contributions.
  \item \emph{Shapley attribution.} The Shapley value of $v$ is
  \[
    \phi_{\mathrm{sup}}=\tau_{\mathrm{sup}}+\tfrac12\iota,\qquad \phi_{\mathrm{dem}}=\tau_{\mathrm{dem}}+\tfrac12\iota,\qquad \phi_{\mathrm{sup}}+\phi_{\mathrm{dem}}=\tau .
  \]
  It is the average of the two path attributions in (a), and it is the only allocation that is efficient, symmetric and additive and gives nothing to a null player. If all four cells are positive, the Shapley value of the game $\log\bar\mu(a_J;b_J)-\log\bar\mu(\mathrm{ref};\mathrm{ref})$ gives the exact factorisation $\mathrm{TR}=\Phi_{\mathrm{sup}}\Phi_{\mathrm{dem}}$, with
  \[
    \Phi_{\mathrm{sup}}=\{\mathrm{AR}(F_{\mathrm{ref}})\,\mathrm{AR}(F_{\mathrm{alg}})\}^{1/2},\qquad
    \Phi_{\mathrm{dem}}=\{\mathrm{DR}(\mathrm{ref})\,\mathrm{DR}(\mathrm{alg})\}^{1/2},
  \]
  and multiplicative interaction $\iota^{\times}=\mathrm{AR}(F_{\mathrm{alg}})/\mathrm{AR}(F_{\mathrm{ref}})=\mathrm{DR}(\mathrm{alg})/\mathrm{DR}(\mathrm{ref})$.
  \item \emph{Bilinear form.} Under \cref{asm:demand}, let $\Delta\mu=\mu_{\mathrm{alg}}-\mu_{\mathrm{ref}}$, $\mu_{\mathrm{mid}}=\frac12(\mu_{\mathrm{alg}}+\mu_{\mathrm{ref}})$ and $F_{\mathrm{mid}}=\frac12(F_{\mathrm{alg}}+F_{\mathrm{ref}})$. Then
  \[
  \begin{aligned}
    \tau_{\mathrm{sup}}&=\mathrm{AA}_{\mathrm{ref}}(F_{\mathrm{ref}}), &
    \tau_{\mathrm{sup}}+\iota&=\mathrm{AA}_{\mathrm{ref}}(F_{\mathrm{alg}}), &
    \phi_{\mathrm{sup}}&=\mathrm{AA}_{\mathrm{ref}}(F_{\mathrm{mid}}),\\
    \tau_{\mathrm{dem}}&={\textstyle\int}\mu_{\mathrm{ref}}\,d(F_{\mathrm{alg}}-F_{\mathrm{ref}}), &
    \iota&={\textstyle\int}\Delta\mu\,d(F_{\mathrm{alg}}-F_{\mathrm{ref}}), &
    \phi_{\mathrm{dem}}&={\textstyle\int}\mu_{\mathrm{mid}}\,d(F_{\mathrm{alg}}-F_{\mathrm{ref}}).
  \end{aligned}
  \]
  \item \emph{When the interaction vanishes.} Under \cref{asm:demand}, $\iota=0$ if and only if $\mathrm{AA}_{\mathrm{ref}}(F_{\mathrm{alg}})=\mathrm{AA}_{\mathrm{ref}}(F_{\mathrm{ref}})$, and $\iota^{\times}=1$ if and only if $\mathrm{AR}_{\mathrm{ref}}(F_{\mathrm{alg}})=\mathrm{AR}_{\mathrm{ref}}(F_{\mathrm{ref}})$. In particular, $\iota=0$ in either of two cases.
  \begin{enumerate}[(i)]
    \item \emph{No demand shaping}: $F_{\mathrm{alg}}=F_{\mathrm{ref}}=F$. Then also $\tau_{\mathrm{dem}}=0$, the two paths in (a) coincide, and observed exposure splits uniquely into a demand-driven part and an algorithm-driven part, $\bar\mu(\mathrm{alg};\mathrm{alg})=\mu_{\mathrm{ref}}(F)+\mathrm{AA}_{\mathrm{ref}}(F)$.
    \item \emph{Homogeneous supply effect}: $\Delta\mu$ is constant on a set $S$ with $F_{\mathrm{alg}}(S)=F_{\mathrm{ref}}(S)=1$.
  \end{enumerate}
  In general, for any such $S$,
  \[
    |\iota|\le\osc_S(\Delta\mu)\,\dTV(F_{\mathrm{alg}},F_{\mathrm{ref}}),\qquad
    |\tau_{\mathrm{dem}}|\le\osc_S(\mu_{\mathrm{ref}})\,\dTV(F_{\mathrm{alg}},F_{\mathrm{ref}}).
  \]
  If $\mathcal D$ is a metric space and $\Delta\mu$ and $\mu_{\mathrm{ref}}$ are Lipschitz with constants $L_\Delta$ and $L_{\mathrm{ref}}$, then $|\iota|\le L_\Delta W_1(F_{\mathrm{alg}},F_{\mathrm{ref}})$ and $|\tau_{\mathrm{dem}}|\le L_{\mathrm{ref}}W_1(F_{\mathrm{alg}},F_{\mathrm{ref}})$.
\end{enumerate}
\end{theorem}

Part (a) is why decompositions of observed exposure are path dependent. A familiar reading of observed exposure takes users as they are and asks what the reference would deliver to them: observed exposure equals $\bar\mu(\mathrm{ref};\mathrm{alg})$ plus the algorithm's increment. That is the demand-first path, and it assigns the whole interaction to supply. A reading that starts from a world without the algorithm, and switches the policy first, assigns it to demand. Neither is wrong. They answer different counterfactual questions, and they disagree by exactly $\iota$.

Part (b) resolves the ambiguity by symmetry. The Shapley value splits the interaction equally between the two factors. By part (c), its supply share is the additive amplification at the average of the two demand distributions, and its demand share is the shift in demand valued at the average of the two policies. This is \citeauthor{kitagawa1955}'s (\citeyear{kitagawa1955}) symmetric decomposition of a difference in rates into a rate component and a composition component, with $\mu_a$ in the role of the rates and $F_b$ in the role of the composition. On the ratio scale the supply share is the geometric mean of the amplification ratios at the two demand distributions. The Shapley value is a convention with axiomatic support, not an identification result. When exposure and mediator interact, direct and indirect effects are not separable even in a randomised trial \citep{robins1992}, and no allocation of $\iota$ is ``the true one''. We therefore report $\tau_{\mathrm{sup}}$, $\tau_{\mathrm{dem}}$ and $\iota$ alongside the Shapley shares.

Part (d) says when the ambiguity disappears. The interaction is bounded by the product of two departures: the heterogeneity of the supply effect across demand profiles, and the extent to which the algorithm has moved demand. It is therefore of second order when both are small. It can be large when both are large, that is, when the algorithm amplifies most for the kinds of demand that it has itself produced. That is the signature of a feedback loop.

In the language of mediation analysis, the formation regime acts on exposure directly and through demand. Here $\tau_{\mathrm{sup}}$ and $\tau_{\mathrm{sup}}+\iota$ play the roles of the pure and total direct effects, $\tau_{\mathrm{dem}}$ and $\tau_{\mathrm{dem}}+\iota$ those of the pure and total indirect effects \citep{robins1992,pearl2001}, and $\iota$ that of the mediated interaction of \citet{vanderweele2014}. Natural direct and indirect effects are cross-world quantities. Under \cref{asm:demand}, the off-diagonal cells \eqref{eq:est-bilinear} instead coincide in expectation with the randomised interventional analogues of \citet{vanderweele2014b}, in which the mediator is drawn from its distribution under the other regime. Unlike those analogues, every cell of our schedule can be produced by an experiment.

\subsection{What a fixed-demand audit identifies}\label{sec:est-audit}

An audit whose puppets never update their demand observes, at best, the persona-level means. Suppose the conditions of \cref{sec:identification} hold, so that the audit identifies $\mu_\pi(d)$ for every $\pi\in\Pi$ and every persona $d$ in its design $\mathcal D_A\subseteq\mathcal D$.

\begin{proposition}[Scope of fixed-demand audits]\label{prop:audit-scope}
Suppose that a fixed-demand audit identifies $\mu_\pi(d)$ for $\pi\in\Pi$ and $d\in\mathcal D_A$, and that \cref{asm:demand} holds.
\begin{enumerate}[(a)]
  \item The audit identifies $\mathrm{AR}_\rho(Q)$ and $\mathrm{AA}_\rho(Q)$ for every reference $\rho$ and every $Q$ concentrated on $\mathcal D_A$. If $F_{\mathrm{ref}}$, $F_{\mathrm{alg}}$ or $F_{\mathrm{mid}}$ is known and concentrated on $\mathcal D_A$, then $Q=F_{\mathrm{ref}}$ identifies $\tau_{\mathrm{sup}}$, $Q=F_{\mathrm{alg}}$ identifies $\tau_{\mathrm{sup}}+\iota$, and $Q=F_{\mathrm{mid}}$ identifies $\phi_{\mathrm{sup}}$.
  \item The law of the audit data does not depend on how demand forms, so the audit data do not identify $\tau_{\mathrm{dem}}$, $\iota$, $\tau$ or $\phi_{\mathrm{dem}}$. Even with oracle knowledge of $\mu_{\mathrm{alg}}$ and $\mu_{\mathrm{ref}}$ on all of $\mathcal D$ and of $F_{\mathrm{alg}}$, the identified set for $\tau$ is an interval with endpoints $\bar\mu(\mathrm{alg};\mathrm{alg})-\sup_{\mathcal D}\mu_{\mathrm{ref}}$ and $\bar\mu(\mathrm{alg};\mathrm{alg})-\inf_{\mathcal D}\mu_{\mathrm{ref}}$, of length $\osc_{\mathcal D}(\mu_{\mathrm{ref}})$. The same holds for $\tau_{\mathrm{dem}}$, with $\mu_{\mathrm{ref}}(F_{\mathrm{alg}})$ in place of $\bar\mu(\mathrm{alg};\mathrm{alg})$. With that oracle knowledge, $\tau$ and $\tau_{\mathrm{dem}}$ are identified if and only if $\mu_{\mathrm{ref}}$ is constant on $\mathcal D$, $\iota$ if and only if $\Delta\mu$ is, and $\phi_{\mathrm{dem}}$ if and only if $\mu_{\mathrm{mid}}$ is. An audit knows these functions only on $\mathcal D_A$ and does not know $F_{\mathrm{alg}}$, so its identified sets are wider still.
  \item \emph{Sensitivity.} $\tau-\mathrm{AA}_{\mathrm{ref}}(F_{\mathrm{alg}})=\tau_{\mathrm{dem}}$ and $\phi_{\mathrm{sup}}-\mathrm{AA}_{\mathrm{ref}}(F_{\mathrm{alg}})=-\iota/2$. If $\mathcal D$ is a metric space, $W_1(F_{\mathrm{alg}},F_{\mathrm{ref}})\le\eta$, and $\mu_{\mathrm{ref}}$ and $\Delta\mu$ are Lipschitz with constants $L_{\mathrm{ref}}$ and $L_\Delta$, then $|\tau-\mathrm{AA}_{\mathrm{ref}}(F_{\mathrm{alg}})|\le L_{\mathrm{ref}}\eta$ and $|\phi_{\mathrm{sup}}-\mathrm{AA}_{\mathrm{ref}}(F_{\mathrm{alg}})|\le L_\Delta\eta/2$.
\end{enumerate}
\end{proposition}

A fixed-demand audit therefore identifies supply-side effects at the personas it audits. The population it speaks for is fixed by $Q$. Speaking for present-day users requires personas that represent $F_{\mathrm{alg}}$, which is itself a measurement problem. Speaking for a world without the algorithm requires $F_{\mathrm{ref}}$, which no audit of the present platform observes. Part (b) is a statement about the best case. Even with the persona-level means known for every conceivable demand profile and present-day demand known, information on demand formation is needed to narrow the total effect beyond the full range of reference exposures. An audit knows less than that. Part (c) replaces that information with a sensitivity parameter $\eta$, the distance by which the algorithm has moved demand, which must come from outside the audit, for instance from panel data on real users.

Four designs reach further, and each has a cost.
\begin{enumerate}[(1)]
  \item A \emph{long-running randomised holdout}, in which some users are served by the reference throughout, identifies $\tau$ for real users, interaction included. The reverse-chronological holdout analysed by \citet{huszar2022} is of this kind.
  \item A \emph{randomised switch} of users from the algorithm to the reference for a fixed period, as in the chronological-feed experiment of \citet{guess2023}, starts at the cell $\bar\mu(\mathrm{ref};\mathrm{alg})$. The contrast with users who stay on the algorithm is $-\mathrm{AA}_{\mathrm{ref}}(F_{\mathrm{alg}})$ when the switch begins. If demand re-forms under the new policy, the contrast moves toward $-\tau$ as the switch lengthens. Demand then changes within the evaluation window, which \cref{asm:demand} excludes, so this regime is described by the schedule cells but not by the bilinear form \eqref{eq:est-bilinear}. The full $2\times2$ schedule, and with it the Shapley split, requires switches in both directions, which only a platform can run.
  \item \emph{Replay designs}, such as the counterfactual bots of \citet{hosseinmardi2024}, copy a real user's viewing history into a bot and then let the bot choose items by a fixed rule. This holds the account at a state formed under the algorithm and varies how later items are chosen. Like a fixed-demand audit, it does not vary the formation regime, and so it says nothing by itself about $F_{\mathrm{ref}}$.
  \item \emph{Adaptive personas}, puppets whose scripted demand updates in response to what they are shown, identify the full $2\times2$ table for synthetic users with the scripted dynamics. Transferring the result to real users requires the script to match human demand dynamics. That is an assumption about people, and no audit of the platform can test it.
\end{enumerate}

\paragraph{Persona contrasts} Parts (a)--(c) of \cref{thm:decomp} use nothing about $F_{\mathrm{alg}}$ and $F_{\mathrm{ref}}$ beyond their being distributions on $\mathcal D$. With point masses at a baseline persona $d_0$ and a persona $d$, the bilinear form \eqref{eq:est-bilinear} becomes the $2\times2$ table $\mu_a(d')$, $a\in\{\mathrm{alg},\mathrm{ref}\}$, $d'\in\{d_0,d\}$. The theorem then splits $\mu_{\mathrm{alg}}(d)-\mu_{\mathrm{ref}}(d_0)$ into a supply contrast $\mathrm{AA}_{\mathrm{ref}}(d_0)$, a demand contrast $\mu_{\mathrm{ref}}(d)-\mu_{\mathrm{ref}}(d_0)$ and an interaction $\mathrm{AA}_{\mathrm{ref}}(d)-\mathrm{AA}_{\mathrm{ref}}(d_0)$, with the same Shapley shares. This persona-contrast decomposition is identified by a fixed-demand audit, and it is the version registered in \cref{sec:protocol}. Its demand contrast compares personas chosen by the researcher. It does not measure demand formed by the algorithm, which \cref{prop:audit-scope}(b) places beyond the reach of such audits.

In the Markov model of \cref{sec:markov}, the supply response is $\mu_\pi(d)=e_T\{P_\pi(d),\alpha_\pi(d)\}$, the expected exposure of a chain with transition matrix $P_\pi(d)$ and initial distribution $\alpha_\pi(d)$. Supply acts through the transition matrix at fixed $d$, and demand shaping moves $d$. The sensitivity results of \cref{sec:markov-perturb} describe the supply side, and they show why a small change in $P$ can produce a large change in exposure.
